% Source: https://artofproblemsolving.com/community/c3h2514203_maximum_and_minimum_of_eulerrsquos_totient_function
% Stable identity: aops-topic-2514203-q01
We use the product formula $\varphi(x)=x\prod_{p\mid x}\left(1-\dfrac1p\right)$, the product being over the distinct primes $p$ dividing $x$ (an empty product is $1$). Hence
\[\dfrac{x}{\varphi(x)}=\prod_{p\mid x}\dfrac{p}{p-1}.\]

\textbf{(a)} Every factor $\dfrac{p}{p-1}$ is greater than $1$. So $\dfrac{x}{\varphi(x)}\ge1$, with equality exactly when $x$ has no prime factor, i.e. $x=1$. The smallest value is $\boxed{1}$, attained only at $x=1$. (For $x\ge2$ the value is larger than $1$, although it can be very close to $1$: for a prime $p$ it is $\dfrac{p}{p-1}$.)

\textbf{(b)} Let $n\ge2$ and $x=n!$. The primes dividing $n!$ are exactly the primes $p\le n$, so, summing geometric series,
\[\dfrac{x}{\varphi(x)}=\prod_{p\le n}\dfrac{1}{1-\dfrac1p}=\prod_{p\le n}\left(1+\dfrac1p+\dfrac1{p^2}+\cdots\right)\ge\prod_{p\le n}\left(1+\dfrac1p+\cdots+\dfrac1{p^m}\right),\]
where $m$ is chosen with $2^m>n$. Expanding the last product gives the sum of $\dfrac1k$ over all integers $k$ whose prime factors are at most $n$ and appear with exponents at most $m$; by unique factorisation each such $k$ appears exactly once. Every $k$ with $1\le k\le n$ is of this kind (a prime power $p^e$ dividing $k$ satisfies $2^e\le p^e\le n<2^m$). Therefore
\[\dfrac{x}{\varphi(x)}\ge1+\dfrac12+\dfrac13+\cdots+\dfrac1n.\]
The harmonic sums are unbounded: for each $j\ge1$,
\[\dfrac{1}{2^{j-1}+1}+\cdots+\dfrac{1}{2^j}\ge2^{j-1}\cdot\dfrac{1}{2^j}=\dfrac12,\qquad\text{so}\qquad1+\dfrac12+\cdots+\dfrac{1}{2^j}\ge1+\dfrac j2.\]
Hence for $x=\left(2^j\right)!$ we get $\dfrac{x}{\varphi(x)}\ge1+\dfrac j2$, which exceeds any given bound for large $j$. So $\dfrac{x}{\varphi(x)}$ has minimum $1$ and no maximum. $\blacksquare$
